\section{Four-dimensional concentration from a fixed occupation strip}
\label{sec:v38-four-dimensional-concentration}
The joint collision operator contains more local information than its
application to a fixed diffusive occupation test. Its expansion is
four-dimensional, and the compact displacement--roof estimate persists
on a fixed, sufficiently small occupation strip. We exploit that strip
without assuming a gap on the entire occupation torus. The conclusions
are a single-index upper bound at the four-dimensional scale and a
local law for windows containing any diverging number of return indices.

Retain $c,\eta_R,A_{m,R},\Omega_R$ and $\mathcal Q_{R,\omega,v}$ from
Section~\ref{sec:v37-occupation-local}. Throughout this section, $a_R,d_R$
are nonnegative multiplier-bounded endpoint families in the sense stated
there. Write $\alpha_R=\nu(a_R)\nu(d_R)$ and
\[
 \Phi_{m,R}(\omega,v)
 =\ell\bigl(M_{d_R}\mathcal Q_{R,\omega,v}^m(a_R\nu)\bigr),
 \qquad \omega=(u,b)\in\T^2\times\R.
\]
The letter $d_R$ here denotes the terminal function, not the lower-right
entry of $\Omega_R$. The physical integral defining $\Phi$ is exactly
\eqref{eq:v37-exact-joint-pairing}; occupation counts the times $[0,m)$.

\subsection{An integrable four-dimensional strip}
\begin{proposition}[The fixed-strip integral]
\label{prop:v38-strip-integral}
For every $B\ge1$ there is $0<\sigma_B<\pi/4$ such that
\begin{equation}\label{eq:v38-strip-integral}
 \int_{[-\pi,\pi]^2}\int_{-B}^{B}\int_{-\sigma_B}^{\sigma_B}
          |\Phi_{m,R}(u,b,v)|\,dv\,db\,du\le C_Bm^{-2}
 \qquad(m\ge1),
\end{equation}
uniformly in $R$ and in any fixed bounded class of endpoints.
On a fixed neighborhood of the four-dimensional origin the integrand
is the sum of a principal term bounded by $C e^{-c_0m|(u,b,v)|^2}$
and an exponentially small complementary term. On the remainder of
the displayed strip its integral is exponentially small.
The width $\sigma_B$ is chosen after $B$; no positive lower bound for
$\sigma_B$ as $B\to\infty$ is asserted.
\end{proposition}
\begin{proof}
Choose $r>0$ so small that the joint spectral splitting in
Proposition~\ref{prop:v37-joint-small-frequency} holds on
$|(u,b,v)|\le3r$. Uniform positivity of $\Omega_R$ and its cubic
remainder imply, after decreasing $r$,
\[
 |\lambda_R(u,b,v)|^m\le e^{-c_0m|(u,b,v)|^2}
 \quad\text{for real }|(u,b,v)|\le3r.
\]
The principal amplitude is $\alpha_R+O(|(u,b,v)|)$, uniformly for
bounded endpoints; the complementary pairing is $O(\rho^m)$.
Choose $\sigma_B<\min(r,\pi/4)$ and also smaller than the occupation
radius supplied by \eqref{eq:v37-compact-small-occupation} for
$|\omega|\ge r$ and $|b|\le B$. The region $|\omega|<r$ is covered
by the joint expansion, while its complement is covered by that
compact estimate. The integral of the Gaussian in four dimensions is
$O(m^{-2})$; the remaining fixed-volume terms are exponentially small.
The norms of the endpoint density and multiplier bound their pairings.
Increasing $C_B$ handles the bounded initial range of $m$.
\end{proof}

For $k\in\Z^2$, $t,\ell\in\R$ put
\begin{equation}\label{eq:v38-central-coordinates}
 z=\left(\frac{k}{\sqrt m},\frac{t-m\bar\tau_R}{\sqrt m}\right),
 \qquad y=\frac{\ell-mc}{\sqrt m},\qquad Z=(z,y).
\end{equation}
For integrable $q,p$ with integrable compactly supported Fourier
transforms define
\[
 I_{m,R}(q,p;k,t,\ell)=\int a_R(x)d_R(T_R^mx)
 \one_{\{K^{\rm c}_{m,R}=k\}}q(S_{m,R}-t)p(A_{m,R}-\ell)\,d\nu(x).
\]
If their Fourier supports lie in $[-B,B]$ and
$[-\sigma_B,\sigma_B]$, respectively, then the exact formula is
\begin{align}\label{eq:v38-double-test-inverse}
 I_{m,R}(q,p;k,t,\ell)
 =\frac1{(2\pi)^4}\int
 &e^{-i[u\cdot k+b(t-m\bar\tau_R)+v(\ell-mc)]}
 \widehat q(-b)\widehat p(-v)\Phi_{m,R}(u,b,v)\,du\,db\,dv.
\end{align}
The integral ranges over that strip. Both minus signs follow from the
convention $\widehat q(b)=\int e^{ibs}q(s)\,ds$. This is Fourier
inversion of two real tests; it is not full occupation-torus inversion.
Fubini is justified by the bounded physical factors and the integrable
Fourier transforms. Proposition~\ref{prop:v38-strip-integral} gives
\begin{equation}\label{eq:v38-double-test-bound}
 |I_{m,R}(q,p;k,t,\ell)|
       \le C_Bm^{-2}\|q\|_1\|p\|_1,
\end{equation}
independently of all targets.

\subsection{Two envelopes, with an order-preserving lower bound}
Use the interval envelopes of
Lemma~\ref{lem:collision-interval-envelopes}. For an interval $J$ and
band $B$ denote them by $q^\pm$, and for $[-H/2,H/2]$ and band
$\sigma$ denote them by $p_H^\pm$. Put
$e_B=2\pi/B$ and $e_\sigma=2\pi/\sigma$. In addition to their
pointwise order and masses, they satisfy
\begin{equation}\label{eq:v38-envelope-L1}
 \|q^\pm-\one_J\|_1=e_B,\quad
 \|p_H^\pm-\one_{[-H/2,H/2]}\|_1=e_\sigma,\quad
 \|p_H^\pm\|_1\le H+e_\sigma.
\end{equation}
Indeed each difference has one sign, and its integral is the prescribed
mass error. In particular
\begin{equation}\label{eq:v38-envelope-Fourier}
 \sup_v\left|\widehat p_H^\pm(v)
       -\widehat{\one}_{[-H/2,H/2]}(v)\right|\le e_\sigma.
\end{equation}
These are the standard one-dimensional Beurling--Selberg envelopes
\cite[Introduction]{CVExt}; no multidimensional extremizer is assumed.

\begin{lemma}[Product envelopes]
\label{lem:v38-product-envelopes}
If $f,g$ are interval indicators and $q^-\le f\le q^+$,
$p^-\le g\le p^+$ are their envelopes, then
\begin{equation}\label{eq:v38-product-envelope}
 q^-p^++q^+p^- -q^+p^+\le fg\le q^+p^+.
\end{equation}
The inequalities remain valid after multiplication by the nonnegative
physical endpoint weight. They do not require either minorant to be
nonnegative.
\end{lemma}
\begin{proof}
The majorants are nonnegative. For the lower bound observe pointwise
that
\[
 q^+p^+-fg=(q^+-f)p^++f(p^+-g)
       \le(q^+-q^-)p^++q^+(p^+-p^-).
\]
Rearranging proves the result. Multiplying two possibly negative
minorants instead would not prove the required lower bound.
\end{proof}

\begin{theorem}[Single-occupation local concentration]
\label{thm:v38-singleton-concentration}
For every bounded roof interval $J$, uniformly for all $R,k,t$, all
integers $n$, and $m\ge1$,
\begin{equation}\label{eq:v38-singleton-concentration}
 \int a_R(x)d_R(T_R^mx)
 \one_{\{K^{\rm c}_{m,R}=k,\ A_{m,R}=n,\ S_{m,R}-t\in J\}}\,d\nu(x)
       \le C(1+|J|)m^{-2}.
\end{equation}
The constant is independent of the position of $J$ and depends only
on the bounded endpoint class and the physical family.
\end{theorem}
\begin{proof}
First take an interval of length at most one. Fix $B=1$ and the strip
width in Proposition~\ref{prop:v38-strip-integral}. A translate of a
single roof majorant of length one bounds its indicator. Because
$A_{m,R}$ is integer-valued,
\[
 \one_{\{A_{m,R}=n\}}=\one_{[-1/2,1/2]}(A_{m,R}-n)
       \le p_1^+(A_{m,R}-n).
\]
The product of these two nonnegative majorants bounds the integrand.
Its integral is bounded by \eqref{eq:v38-double-test-bound}, with
$\|q^+\|_1=1+2\pi$ and $\|p_1^+\|_1=1+2\pi/\sigma_1$.
Translation changes no Fourier absolute bound. This proves the result
for every interval of length at most one, including any choice of
endpoints. Cover an arbitrary bounded interval by at most $|J|+2$
such intervals and sum. No asymptotic lower bound for a singleton is
used or obtained here.
\end{proof}

\subsection{Moving tests at an arbitrarily slow diverging width}
For $H\ge1$ set
\begin{equation}\label{eq:v38-Gaussian-window}
 \mathcal G_{m,H,R}(z,y)=\frac1H\int_{-H/2}^{H/2}
           g_{\Omega_R}(z,y+s/\sqrt m)\,ds.
\end{equation}
This expression retains the variation of the Gaussian across the
whole window, including windows of diffusive size.

\begin{lemma}[A uniform moving-envelope comparison]
\label{lem:v38-moving-envelope}
Fix $B$, $0<\sigma\le\sigma_B$, and a roof test $q$ whose integrable
Fourier transform is supported in $[-B,B]$. For each $H\ge1$ let
$p_H^\pm$ be the interval envelopes with band $\sigma$. There are
$\epsilon_m\to0$ and $C_q<\infty$, depending on the fixed bands and
$q$ but not on $H,R,k,t,\ell$, such that
\begin{align}\label{eq:v38-moving-envelope}
 \left|\frac{m^2}{H}I_{m,R}(q,p_H^\pm;k,t,\ell)
       -\alpha_R\left(\int q\right)\mathcal G_{m,H,R}(z,y)\right|
       \le \epsilon_m+\frac{C_q}{\sigma H}.
\end{align}
It suffices here that $q$ is fixed; no convergence rate for
$\epsilon_m$ is asserted.
\end{lemma}
\begin{proof}
Apply \eqref{eq:v38-double-test-inverse}. Divide its Fourier weight by
$H$ and use $|\widehat p_H^\pm|/H\le1+e_\sigma/H$.
Outside the fixed origin neighborhood, the scaled inverse is at most
$C_Bm^2\rho_B^m(1+e_\sigma/H)$ and tends to zero uniformly for
$H\ge1$. Inside it put $w=\sqrt m(u,b,v)$. The spectral expansion,
its Gaussian bound and the principal amplitude give an $L^1(dw)$
error tending uniformly to zero when the operator pairing is replaced
by $\alpha_Re^{-w^{\mathsf T}\Omega_Rw/2}$. The same holds when
$\widehat q(-w_3/\sqrt m)$ is replaced by $\int q$. The bounded
factor $1+e_\sigma/H$ is at most $1+e_\sigma$; thus all these errors
are absorbed in one $\epsilon_m$ independent of $H$. The inverse
oscillation has modulus one, so this reasoning is uniform in all
central coordinates, not just convergent target sequences.

Replace $\widehat p_H^\pm(-w_4/\sqrt m)/H$ by
$\widehat{\one}_{[-H/2,H/2]}(-w_4/\sqrt m)/H$.
Equation~\eqref{eq:v38-envelope-Fourier} and the integrable Gaussian
bound make the error at most $C_q/(\sigma H)$. Extending the remaining
Gaussian integral to $\R^4$ costs another uniformly vanishing error.
Finally, Fourier inversion of the Gaussian, followed by integration
of the interval indicator, gives exactly
\[
 \frac1{(2\pi)^4}\int_{\R^4}e^{-iw\cdot(z,y)}
 e^{-w^{\mathsf T}\Omega_Rw/2}
 \frac1H\int_{-H/2}^{H/2}e^{-iw_4s/\sqrt m}\,ds\,dw
       =\mathcal G_{m,H,R}(z,y).
\]
The sign is consistent with testing $A_m-\ell$. The centered interval
is symmetric, but no sign is discarded in this calculation.
\end{proof}

\begin{theorem}[Local law in an arbitrarily slowly widening occupation block]
\label{thm:v38-occupation-mesoscopic}
Let $J$ be a fixed bounded interval of positive length. Let
$\mathcal B=\{j,j+1,\ldots,j+H-1\}$ be a block of $H\ge1$ integers
with center $\ell=j+(H-1)/2$, and use \eqref{eq:v38-central-coordinates}.
For every sequence $h_m\to\infty$, uniformly over $H\ge h_m$, $R$, and
$|Z|\le M$,
\begin{align}\label{eq:v38-occupation-mesoscopic}
 &\frac{m^2}{H}\int a_R(x)d_R(T_R^mx)
 \one_{\{K^{\rm c}_{m,R}=k,\ A_{m,R}\in\mathcal B,\ S_{m,R}-t\in J\}}
          \,d\nu(x)
 -\alpha_R|J|\mathcal G_{m,H,R}(z,y)\longrightarrow0.
\end{align}
In particular, if also $H\le u_m\sqrt m$ with $u_m\to0$, the Gaussian
average can be replaced by $g_{\Omega_R}(z,y)$. There is no lower
polynomial bound on $h_m$: every diverging width is admissible.
\end{theorem}
\begin{proof}
The half-integer endpoints of the interval
$[j-1/2,j+H-1/2]$ select exactly the block $\mathcal B$, so its
indicator is the centered interval test of length $H$ at $\ell$.
First fix $B\ge1$, then fix $\sigma\le\sigma_B$ and choose the two
pairs of envelopes. Apply Lemma~\ref{lem:v38-product-envelopes}.
Lemma~\ref{lem:v38-moving-envelope} evaluates all three terms in the
lower bound and the one term in the upper bound. In the main term the
upper roof mass is $|J|+e_B$; the lower combination has roof mass
$(|J|-e_B)+(|J|+e_B)-(|J|+e_B)=|J|-e_B$.
Uniform boundedness of $\alpha_R$ and $g_{\Omega_R}$ therefore gives
\begin{equation}\label{eq:v38-two-band-budget}
 \left|\frac{m^2}{H}\int a_R(d_R\circ T_R^m)
       \one_{\{K_m^{\rm c}=k,A_m\in\mathcal B,S_m-t\in J\}}\,d\nu
       -\alpha_R|J|\mathcal G_{m,H,R}(z,y)\right|
 \le \epsilon_m(B)+\frac{C_J}{B}+\frac{C_J}{\sigma H},
\end{equation}
where $\epsilon_m(B)\to0$ for the already fixed $B,\sigma$.
Take $m\to\infty$ uniformly for $H\ge h_m$ and only afterwards
$B\to\infty$. The width $\sigma$ may become smaller when $B$ changes;
this causes no problem because it was fixed before the first limit.
No $m$-dependent band is inserted into the spectral theorem.

When $H/\sqrt m\to0$, uniform continuity of the Gaussian densities
on central compact sets makes the average in
\eqref{eq:v38-Gaussian-window} differ from $g_{\Omega_R}(z,y)$ by
$o(1)$ uniformly. This proves the final assertion. The argument is a
uniform moving-test estimate; weak convergence for fixed diffusive
intervals alone would not justify it.
\end{proof}
